% generated by paper_nd.py
\begin{table}[t]
\centering
\begin{tabular}{ll|cc|cc|cc}
\hline\hline
& & \multicolumn{2}{c|}{$N_d=3$} & \multicolumn{2}{c|}{$N_d=5$} & \multicolumn{2}{c}{$N_d=7$}\\
$N_d\alpha_d$ & $\alpha_d\,(N_d=3)$ & LO & NLO & LO & NLO & LO & NLO\\
\hline
0.003 & 0.001 & 33.53 & 33.52 & 48.8(5) & 48.4(5) & 57.1(6) & 56.4(6)\\
0.009 & 0.003 & 32.15 & 32.13 & 46.6(5) & 46.1(5) & 53.8(6) & 53.9(6)\\
0.03 & 0.01 & 28.47 & 28.39 & 40.9(5) & 41.0(5) & 47.9(5) & 47.3(5)\\
0.06 & 0.02 & 24.56 & 24.38 & 34.6(4) & 34.2(3) & 40.2(4) & 40.4(4)\\
0.09 & 0.03 & 21.50 & 21.21 & 30.6(3) & 30.3(3) & 35.6(4) & 34.3(3)\\
0.15 & 0.05 & 16.81 & 16.38 & 23.6(2) & 23.3(2) & 27.3(3) & 26.4(2)\\
0.21 & 0.07 & 13.32 & 12.80 & 18.7(2) & 17.9(2) & 21.7(2) & 20.4(2)\\
\hline\hline
\end{tabular}
\caption{At the same $N_d\alpha_d(m_\psi)$ the dark baryon is larger for larger $N_d$, by a factor that depends little on the coupling. We give $\langle r^2\rangle m_\psi^2(N_d\alpha_d)^2$ at $\lambda=1$. The $N_d=3$ values are the variational solutions of Table~\ref{tab:smallalpha}, and the $N_d=5,7$ values are GFMC (App.~\ref{app:gfmc:nd}), with the uncertainty in parentheses, including that of the trial correction. Fig.~\ref{fig:nd} shows the variation with $\lambda$.}
\label{tab:nd}
\end{table}
